% Seminar 4 - Page 3. Fragment for
% % Seminar 4 - Page 3. Fragment for
% % Seminar 4 - Page 3. Fragment for
% % Seminar 4 - Page 3. Fragment for
% \input{seminars/pages/seminar4-pagina3-en}; no preamble or forced page break.
% EDITORIAL NOTES:
% - In Exercise 7 the manuscript appears to say "nu e inj." (not
%   injective). The canonical inclusion Z -> Q IS injective, but not
%   surjective. This apparent slip is retained in the displayed text,
%   marked "as written", rather than silently replaced.
% - All monoids in Exercise 7 are considered under multiplication.
%   The diagram and the argument use the epimorphism definition: two
%   morphisms out of Q agreeing on Z must be equal.
% - The proof writes e=u(1)=u(m/m) and similarly for v, then uses
%   uniqueness of inverses in a monoid. The handwritten line for v
%   uses a long dash in place of repeated intermediate steps.
% - Exercise 8 assumes a hypothetical isomorphism (not merely an
%   arbitrary bijection), as shown by the homomorphism identity and
%   use of injectivity at the end.
% - The handwritten right-hand brace at f(a)f(b)=4 contains a faint
%   second entry f(b)=...?; the decisive f(a)=1 => a=1 is clear.
% - The primes x_1,x_2,x_3 are prime in N^*, not "prime" elements of
%   3N+1. The word "prime" is taken verbatim from the manuscript.
% - Faint show-through and crossed-out intermediate notes are omitted.

\subsection*{Seminar 4 - Page 3}

\noindent\underline{Exercise 7.}
The inclusion
\[
 i:(\mathbb Z,\cdot)\longrightarrow(\mathbb Q,\cdot),
 \qquad i(x)=x,
\]
is an epimorphism of monoids which is \emph{not injective}
(as written in the manuscript).

\smallskip
\noindent\underline{Epimorphism.}
For every monoid $M$ and all monoid homomorphisms $u,v$ such that
$u\circ i=v\circ i$, one has $u=v$.
\[
\begin{tikzcd}[column sep=large]
 \mathbb Z \arrow[r,"i"] &
 \mathbb Q \arrow[r,shift left=.7ex,"u"]
            \arrow[r,shift right=.7ex,swap,"v"] & M
\end{tikzcd}
\]

\noindent Let $(M,\cdot)$ be a monoid and
$u,v:\mathbb Q\to M$ monoid homomorphisms with
$u\circ i=v\circ i$. Hence
\[
 u(m)=v(m)\qquad(\forall m\in\mathbb Z).
\]
We wish to show $u(r)=v(r)$ for every $r\in\mathbb Q$.
Let $m\in\mathbb Z$, $m\ne0$. Then
\[
\begin{aligned}
 e=u(1)
  &=u\!\left(\frac{m}{m}\right)
   =u\!\left(m\cdot\frac1m\right)
   =u(m)u\!\left(\frac1m\right)
   =u\!\left(\frac1m\right)u(m),\\
 e=v(1)
  &=\cdots
   =v\!\left(\frac1m\right)v(m).
\end{aligned}
\]
Therefore $u(1/m)$ is the inverse of $u(m)$, and similarly
$v(1/m)$ is the inverse of $v(m)$. Since $u(m)=v(m)$,
\[
 u\!\left(\frac1m\right)=v\!\left(\frac1m\right).
\]
For $r=n/m\in\mathbb Q$, with $m\ne0$,
\[
 u(r)=u(n)u\!\left(\frac1m\right)
     =v(n)v\!\left(\frac1m\right)=v(r).
\]
Consequently $u=v$.

\smallskip
\noindent [Every epimorphism of groups is surjective.]

\medskip
\noindent\underline{Exercise 8.}
The monoids
\[
 (\mathbb N^*,\cdot),\qquad
 (3\mathbb N+1,\cdot),\qquad
 3\mathbb N+1=\{3k+1\mid k\in\mathbb N\},
\]
are not isomorphic.

\noindent\underline{Proof.}
Suppose that there exists an isomorphism
\[
 f:\mathbb N^*\longrightarrow 3\mathbb N+1,
 \qquad f(xy)=f(x)f(y)
 \quad(\forall x,y\in\mathbb N^*).
\]
Since $f$ is surjective, there exist $x_1,x_2,x_3\in\mathbb N^*$
such that
\[
 f(x_1)=4,\qquad f(x_2)=10,\qquad f(x_3)=25.
\]
We show that $x_1,x_2,x_3$ are prime numbers.
If $x_1=ab$ with $a,b>1$, then
\[
 f(x_1)=f(a)f(b)=4
 \quad\Longrightarrow\quad f(a)=1
 \quad\Longrightarrow\quad a=1,
 \qquad\text{a contradiction.}
\]
Analogously, $x_2$ and $x_3$ are prime numbers.
But $4\cdot25=10^2$, so
\[
\begin{aligned}
 f(x_1)f(x_3)&=f(x_2)^2,\\
 f(x_1x_3)&=f(x_2^2),\\
 x_1x_3&=x_2^2
  \quad\Longrightarrow\quad x_1=x_2=x_3,
  \qquad\text{false.}
\end{aligned}
\]
; no preamble or forced page break.
% EDITORIAL NOTES:
% - In Exercise 7 the manuscript appears to say "nu e inj." (not
%   injective). The canonical inclusion Z -> Q IS injective, but not
%   surjective. This apparent slip is retained in the displayed text,
%   marked "as written", rather than silently replaced.
% - All monoids in Exercise 7 are considered under multiplication.
%   The diagram and the argument use the epimorphism definition: two
%   morphisms out of Q agreeing on Z must be equal.
% - The proof writes e=u(1)=u(m/m) and similarly for v, then uses
%   uniqueness of inverses in a monoid. The handwritten line for v
%   uses a long dash in place of repeated intermediate steps.
% - Exercise 8 assumes a hypothetical isomorphism (not merely an
%   arbitrary bijection), as shown by the homomorphism identity and
%   use of injectivity at the end.
% - The handwritten right-hand brace at f(a)f(b)=4 contains a faint
%   second entry f(b)=...?; the decisive f(a)=1 => a=1 is clear.
% - The primes x_1,x_2,x_3 are prime in N^*, not "prime" elements of
%   3N+1. The word "prime" is taken verbatim from the manuscript.
% - Faint show-through and crossed-out intermediate notes are omitted.

\subsection*{Seminar 4 - Page 3}

\noindent\underline{Exercise 7.}
The inclusion
\[
 i:(\mathbb Z,\cdot)\longrightarrow(\mathbb Q,\cdot),
 \qquad i(x)=x,
\]
is an epimorphism of monoids which is \emph{not injective}
(as written in the manuscript).

\smallskip
\noindent\underline{Epimorphism.}
For every monoid $M$ and all monoid homomorphisms $u,v$ such that
$u\circ i=v\circ i$, one has $u=v$.
\[
\begin{tikzcd}[column sep=large]
 \mathbb Z \arrow[r,"i"] &
 \mathbb Q \arrow[r,shift left=.7ex,"u"]
            \arrow[r,shift right=.7ex,swap,"v"] & M
\end{tikzcd}
\]

\noindent Let $(M,\cdot)$ be a monoid and
$u,v:\mathbb Q\to M$ monoid homomorphisms with
$u\circ i=v\circ i$. Hence
\[
 u(m)=v(m)\qquad(\forall m\in\mathbb Z).
\]
We wish to show $u(r)=v(r)$ for every $r\in\mathbb Q$.
Let $m\in\mathbb Z$, $m\ne0$. Then
\[
\begin{aligned}
 e=u(1)
  &=u\!\left(\frac{m}{m}\right)
   =u\!\left(m\cdot\frac1m\right)
   =u(m)u\!\left(\frac1m\right)
   =u\!\left(\frac1m\right)u(m),\\
 e=v(1)
  &=\cdots
   =v\!\left(\frac1m\right)v(m).
\end{aligned}
\]
Therefore $u(1/m)$ is the inverse of $u(m)$, and similarly
$v(1/m)$ is the inverse of $v(m)$. Since $u(m)=v(m)$,
\[
 u\!\left(\frac1m\right)=v\!\left(\frac1m\right).
\]
For $r=n/m\in\mathbb Q$, with $m\ne0$,
\[
 u(r)=u(n)u\!\left(\frac1m\right)
     =v(n)v\!\left(\frac1m\right)=v(r).
\]
Consequently $u=v$.

\smallskip
\noindent [Every epimorphism of groups is surjective.]

\medskip
\noindent\underline{Exercise 8.}
The monoids
\[
 (\mathbb N^*,\cdot),\qquad
 (3\mathbb N+1,\cdot),\qquad
 3\mathbb N+1=\{3k+1\mid k\in\mathbb N\},
\]
are not isomorphic.

\noindent\underline{Proof.}
Suppose that there exists an isomorphism
\[
 f:\mathbb N^*\longrightarrow 3\mathbb N+1,
 \qquad f(xy)=f(x)f(y)
 \quad(\forall x,y\in\mathbb N^*).
\]
Since $f$ is surjective, there exist $x_1,x_2,x_3\in\mathbb N^*$
such that
\[
 f(x_1)=4,\qquad f(x_2)=10,\qquad f(x_3)=25.
\]
We show that $x_1,x_2,x_3$ are prime numbers.
If $x_1=ab$ with $a,b>1$, then
\[
 f(x_1)=f(a)f(b)=4
 \quad\Longrightarrow\quad f(a)=1
 \quad\Longrightarrow\quad a=1,
 \qquad\text{a contradiction.}
\]
Analogously, $x_2$ and $x_3$ are prime numbers.
But $4\cdot25=10^2$, so
\[
\begin{aligned}
 f(x_1)f(x_3)&=f(x_2)^2,\\
 f(x_1x_3)&=f(x_2^2),\\
 x_1x_3&=x_2^2
  \quad\Longrightarrow\quad x_1=x_2=x_3,
  \qquad\text{false.}
\end{aligned}
\]
; no preamble or forced page break.
% EDITORIAL NOTES:
% - In Exercise 7 the manuscript appears to say "nu e inj." (not
%   injective). The canonical inclusion Z -> Q IS injective, but not
%   surjective. This apparent slip is retained in the displayed text,
%   marked "as written", rather than silently replaced.
% - All monoids in Exercise 7 are considered under multiplication.
%   The diagram and the argument use the epimorphism definition: two
%   morphisms out of Q agreeing on Z must be equal.
% - The proof writes e=u(1)=u(m/m) and similarly for v, then uses
%   uniqueness of inverses in a monoid. The handwritten line for v
%   uses a long dash in place of repeated intermediate steps.
% - Exercise 8 assumes a hypothetical isomorphism (not merely an
%   arbitrary bijection), as shown by the homomorphism identity and
%   use of injectivity at the end.
% - The handwritten right-hand brace at f(a)f(b)=4 contains a faint
%   second entry f(b)=...?; the decisive f(a)=1 => a=1 is clear.
% - The primes x_1,x_2,x_3 are prime in N^*, not "prime" elements of
%   3N+1. The word "prime" is taken verbatim from the manuscript.
% - Faint show-through and crossed-out intermediate notes are omitted.

\subsection*{Seminar 4 - Page 3}

\noindent\underline{Exercise 7.}
The inclusion
\[
 i:(\mathbb Z,\cdot)\longrightarrow(\mathbb Q,\cdot),
 \qquad i(x)=x,
\]
is an epimorphism of monoids which is \emph{not injective}
(as written in the manuscript).

\smallskip
\noindent\underline{Epimorphism.}
For every monoid $M$ and all monoid homomorphisms $u,v$ such that
$u\circ i=v\circ i$, one has $u=v$.
\[
\begin{tikzcd}[column sep=large]
 \mathbb Z \arrow[r,"i"] &
 \mathbb Q \arrow[r,shift left=.7ex,"u"]
            \arrow[r,shift right=.7ex,swap,"v"] & M
\end{tikzcd}
\]

\noindent Let $(M,\cdot)$ be a monoid and
$u,v:\mathbb Q\to M$ monoid homomorphisms with
$u\circ i=v\circ i$. Hence
\[
 u(m)=v(m)\qquad(\forall m\in\mathbb Z).
\]
We wish to show $u(r)=v(r)$ for every $r\in\mathbb Q$.
Let $m\in\mathbb Z$, $m\ne0$. Then
\[
\begin{aligned}
 e=u(1)
  &=u\!\left(\frac{m}{m}\right)
   =u\!\left(m\cdot\frac1m\right)
   =u(m)u\!\left(\frac1m\right)
   =u\!\left(\frac1m\right)u(m),\\
 e=v(1)
  &=\cdots
   =v\!\left(\frac1m\right)v(m).
\end{aligned}
\]
Therefore $u(1/m)$ is the inverse of $u(m)$, and similarly
$v(1/m)$ is the inverse of $v(m)$. Since $u(m)=v(m)$,
\[
 u\!\left(\frac1m\right)=v\!\left(\frac1m\right).
\]
For $r=n/m\in\mathbb Q$, with $m\ne0$,
\[
 u(r)=u(n)u\!\left(\frac1m\right)
     =v(n)v\!\left(\frac1m\right)=v(r).
\]
Consequently $u=v$.

\smallskip
\noindent [Every epimorphism of groups is surjective.]

\medskip
\noindent\underline{Exercise 8.}
The monoids
\[
 (\mathbb N^*,\cdot),\qquad
 (3\mathbb N+1,\cdot),\qquad
 3\mathbb N+1=\{3k+1\mid k\in\mathbb N\},
\]
are not isomorphic.

\noindent\underline{Proof.}
Suppose that there exists an isomorphism
\[
 f:\mathbb N^*\longrightarrow 3\mathbb N+1,
 \qquad f(xy)=f(x)f(y)
 \quad(\forall x,y\in\mathbb N^*).
\]
Since $f$ is surjective, there exist $x_1,x_2,x_3\in\mathbb N^*$
such that
\[
 f(x_1)=4,\qquad f(x_2)=10,\qquad f(x_3)=25.
\]
We show that $x_1,x_2,x_3$ are prime numbers.
If $x_1=ab$ with $a,b>1$, then
\[
 f(x_1)=f(a)f(b)=4
 \quad\Longrightarrow\quad f(a)=1
 \quad\Longrightarrow\quad a=1,
 \qquad\text{a contradiction.}
\]
Analogously, $x_2$ and $x_3$ are prime numbers.
But $4\cdot25=10^2$, so
\[
\begin{aligned}
 f(x_1)f(x_3)&=f(x_2)^2,\\
 f(x_1x_3)&=f(x_2^2),\\
 x_1x_3&=x_2^2
  \quad\Longrightarrow\quad x_1=x_2=x_3,
  \qquad\text{false.}
\end{aligned}
\]
; no preamble or forced page break.
% EDITORIAL NOTES:
% - In Exercise 7 the manuscript appears to say "nu e inj." (not
%   injective). The canonical inclusion Z -> Q IS injective, but not
%   surjective. This apparent slip is retained in the displayed text,
%   marked "as written", rather than silently replaced.
% - All monoids in Exercise 7 are considered under multiplication.
%   The diagram and the argument use the epimorphism definition: two
%   morphisms out of Q agreeing on Z must be equal.
% - The proof writes e=u(1)=u(m/m) and similarly for v, then uses
%   uniqueness of inverses in a monoid. The handwritten line for v
%   uses a long dash in place of repeated intermediate steps.
% - Exercise 8 assumes a hypothetical isomorphism (not merely an
%   arbitrary bijection), as shown by the homomorphism identity and
%   use of injectivity at the end.
% - The handwritten right-hand brace at f(a)f(b)=4 contains a faint
%   second entry f(b)=...?; the decisive f(a)=1 => a=1 is clear.
% - The primes x_1,x_2,x_3 are prime in N^*, not "prime" elements of
%   3N+1. The word "prime" is taken verbatim from the manuscript.
% - Faint show-through and crossed-out intermediate notes are omitted.

\subsection*{Seminar 4 - Page 3}

\noindent\underline{Exercise 7.}
The inclusion
\[
 i:(\mathbb Z,\cdot)\longrightarrow(\mathbb Q,\cdot),
 \qquad i(x)=x,
\]
is an epimorphism of monoids which is \emph{not injective}
(as written in the manuscript).

\smallskip
\noindent\underline{Epimorphism.}
For every monoid $M$ and all monoid homomorphisms $u,v$ such that
$u\circ i=v\circ i$, one has $u=v$.
\[
\begin{tikzcd}[column sep=large]
 \mathbb Z \arrow[r,"i"] &
 \mathbb Q \arrow[r,shift left=.7ex,"u"]
            \arrow[r,shift right=.7ex,swap,"v"] & M
\end{tikzcd}
\]

\noindent Let $(M,\cdot)$ be a monoid and
$u,v:\mathbb Q\to M$ monoid homomorphisms with
$u\circ i=v\circ i$. Hence
\[
 u(m)=v(m)\qquad(\forall m\in\mathbb Z).
\]
We wish to show $u(r)=v(r)$ for every $r\in\mathbb Q$.
Let $m\in\mathbb Z$, $m\ne0$. Then
\[
\begin{aligned}
 e=u(1)
  &=u\!\left(\frac{m}{m}\right)
   =u\!\left(m\cdot\frac1m\right)
   =u(m)u\!\left(\frac1m\right)
   =u\!\left(\frac1m\right)u(m),\\
 e=v(1)
  &=\cdots
   =v\!\left(\frac1m\right)v(m).
\end{aligned}
\]
Therefore $u(1/m)$ is the inverse of $u(m)$, and similarly
$v(1/m)$ is the inverse of $v(m)$. Since $u(m)=v(m)$,
\[
 u\!\left(\frac1m\right)=v\!\left(\frac1m\right).
\]
For $r=n/m\in\mathbb Q$, with $m\ne0$,
\[
 u(r)=u(n)u\!\left(\frac1m\right)
     =v(n)v\!\left(\frac1m\right)=v(r).
\]
Consequently $u=v$.

\smallskip
\noindent [Every epimorphism of groups is surjective.]

\medskip
\noindent\underline{Exercise 8.}
The monoids
\[
 (\mathbb N^*,\cdot),\qquad
 (3\mathbb N+1,\cdot),\qquad
 3\mathbb N+1=\{3k+1\mid k\in\mathbb N\},
\]
are not isomorphic.

\noindent\underline{Proof.}
Suppose that there exists an isomorphism
\[
 f:\mathbb N^*\longrightarrow 3\mathbb N+1,
 \qquad f(xy)=f(x)f(y)
 \quad(\forall x,y\in\mathbb N^*).
\]
Since $f$ is surjective, there exist $x_1,x_2,x_3\in\mathbb N^*$
such that
\[
 f(x_1)=4,\qquad f(x_2)=10,\qquad f(x_3)=25.
\]
We show that $x_1,x_2,x_3$ are prime numbers.
If $x_1=ab$ with $a,b>1$, then
\[
 f(x_1)=f(a)f(b)=4
 \quad\Longrightarrow\quad f(a)=1
 \quad\Longrightarrow\quad a=1,
 \qquad\text{a contradiction.}
\]
Analogously, $x_2$ and $x_3$ are prime numbers.
But $4\cdot25=10^2$, so
\[
\begin{aligned}
 f(x_1)f(x_3)&=f(x_2)^2,\\
 f(x_1x_3)&=f(x_2^2),\\
 x_1x_3&=x_2^2
  \quad\Longrightarrow\quad x_1=x_2=x_3,
  \qquad\text{false.}
\end{aligned}
\]
